% V.2c and V.3: this chapter must cover the background, the objectives and the
% work plan, and together with chapters 6 to 8 it has to reach 30 pages (25,
% plus 5 for the second author).

\chapter{Introduction}
\label{cap:introduction}

\section{Background}
\label{sec:background}

\todo{Musical similarity has no operational definition. State the problem the
way the literature does: recommendation, playlist generation and catalogue
exploration all need a notion of "sounds like", and none of the public corpora
annotates it directly. Explain what is used instead, and why each substitute is
imperfect: genre labels chosen by the artist, tags written by listeners, and the
small sets of human triplet judgements that do exist. Cite \cite{Ellis2002} and
\cite{Logan2003} for the ground-truth problem itself.}

\section{Objectives}
\label{sec:objectives}

\todo{State the two questions the work answers, which are the two levels of the
evaluation:
\begin{itemize}
  \item \textbf{N1.} Does the representation contain musical information?
        Measured with probing classifiers over frozen embeddings, following
        \cite{Castellon2021} and \cite{Yuan2023}.
  \item \textbf{N2.} Does the representation induce a good similarity graph?
        Measured with retrieval, graph structure, robustness, and agreement with
        the human triplet judgements of MagnaTagATune.
\end{itemize}
Add the question that motivates the whole study: does the genre proxy measure
the same thing as perceived similarity?}

\section{Work plan}
\label{sec:work-plan}

\todo{Describe the plan as it was actually followed: a reproducible experiment
framework first, then the MFCC baseline, then the graph and retrieval
evaluation, then the ground truth work, and finally the trained models and MERT
on the department server. Say which parts ran on a laptop and which needed a
GPU. Reference Appendix \ref{Appendix:reproducibility} for the exact commands,
and give the repository URL required by V.10: \repositorio.}
